\section{1. Introduction}
Distributed average-consensus algorithms formalize how decentralized networks of autonomous agents reach shared numerical agreements through local message passing\cite{degroot1974reaching}. In classical frameworks such as the DeGroot model, agents update their scalar estimates over discrete time steps through convex combinations of their neighbors' prior states\cite{degroot1974reaching}. Most extensions to this architecture assume that communication channels are static, unweighted, or adaptively weighted to prioritize accuracy, historical prestige, or metric proximity\cite{degroot1974reaching}. Under standard conditions of persistent connectivity and aperiodicity, these systems contract toward a static consensus lying strictly within the convex hull of initial network estimates\cite{degroot1974reaching}.

Empirical sociology and organizational psychology document a persistent counter-behavior in human collectives: groups frequently penalize their highest-performing members\cite{degroot1974reaching}. This behavioral tendency, formalized in social psychology as "tall poppy syndrome," describes the social derogation, exclusion, or sanctioning of individuals whose achievements or accuracy significantly surpass the group average\cite{degroot1974reaching}. While leveling mechanisms, reverse dominance hierarchies, and anti-outlier pressures are documented across varied institutional and cultural contexts, asymmetric status-based penalties targeting positive performance have been absent from mathematical opinion dynamics\cite{degroot1974reaching}.

This paper addresses that gap by introducing an endogenous, status-based weight suppression operator into a distributed consensus network\cite{degroot1974reaching}. The central question driving this work is operational: if an agent's outgoing broadcast influence is systematically reduced as a function of its positive deviation from the collective mean, does this leveling pressure impede or destabilize network consensus\cite{degroot1974reaching}? In estimation tasks where an outlier pulls the collective toward a ground truth, this mechanism penalizes the network's most informative agent precisely because of its accuracy\cite{degroot1974reaching}.

Analyzing this mechanism requires moving beyond linear systems theory\cite{degroot1974reaching}. Because influence weights vary at each discrete step based on instantaneous agent opinions, the interaction topology forms a closed-loop, state-dependent dynamical system\cite{degroot1974reaching,hajnal1958weak}. Previous treatments of time-varying consensus frequently rely on open-loop non-homogeneous Markov chain ergodicity, using Hajnal contraction coefficients to evaluate whether an exogenous sequence of stochastic matrices contracts to rank one\cite{degroot1974reaching,hajnal1958weak}. In the presence of cognitive anchoring, however, the governing difference equation is affine rather than linear\cite{acemoglu2013opinion,hajnal1958weak}. The network transition matrix remains strictly row-stochastic by construction, ensuring that estimates remain bounded within the convex hull of initial conditions and preventing numerical divergence\cite{degroot1974reaching}. The appropriate analytical framework is therefore the theory of switched affine systems and discrete-time bifurcations\cite{acemoglu2013opinion,hajnal1958weak}. The peer interaction operator acts as a contraction mapping on the disagreement subspace, while cognitive anchors continuously inject initial state dispersion back into the network\cite{acemoglu2013opinion,hajnal1958weak}.

This formulation introduces a cognitive resistance parameter: partial stubbornness ($\gamma_{stub} \in [0, 1)$)\cite{degroot1974reaching}. In a pure averaging update without cognitive memory ($\gamma_{stub} = 0$), an outlier whose outgoing influence is zeroed still absorbs incoming influence from peers\cite{degroot1974reaching}. The outlier inevitably assimilates into the unpenalized group average, erasing its private information and collapsing the network into a biased consensus\cite{degroot1974reaching}. Anchoring agents partially to their initial private values models cognitive dissonance and resistance to group conformity, creating persistent structural tension between group leveling and individual belief retention\cite{degroot1974reaching}.

Computational simulations across 368,991 independent network evaluations spanning Erdős-Rényi, Barabási-Albert, Watts-Strogatz, and regular topologies show that combining asymmetric suppression with cognitive anchoring produces four distinct stability regimes\cite{degroot1974reaching,acemoglu2013opinion}:

Regime 1 (Biased Convergence): Under weak suppression ($P < 0.3$) or absent stubbornness ($\gamma_{stub} = 0$), the network contracts to a uniform consensus\cite{degroot1974reaching,acemoglu2013opinion}. The final consensus value is systematically biased away from the true mean because unpenalized lower-performing agents are over-weighted during intermediate updates\cite{degroot1974reaching}.

Regime 1B (Fixed Dissensus): Under moderate penalties combined with non-zero stubbornness ($\gamma_{stub} > 0$), peer influence weights stabilize, and opinion dynamics freeze into a stationary, non-consensus configuration\cite{acemoglu2013opinion}. Agents maintain persistent, non-zero spatial variance without moving over time\cite{acemoglu2013opinion}.

Regime 2 (Slow Mixing): Near transition boundaries, state trajectories exhibit prolonged relaxation transients\cite{degroot1974reaching,acemoglu2013opinion}. Interaction weights chatter across the state-dependent switching boundary, producing slow algebraic mixing without entering clean periodic orbits\cite{acemoglu2013opinion,hajnal1958weak}.

Regime 3 (Non-convergent Cycling): Under maximal suppression ($P = 1.0$) paired with stubborn anchoring ($\gamma_{stub} > 0$), the system undergoes a bifurcation into a stable, bounded limit cycle\cite{degroot1974reaching,acemoglu2013opinion}. Accurate outliers are silenced, the collective drifts downward, the standardized deviation of the outliers normalizes, their influence is restored, they pull the group upward, and the penalty reactivates\cite{degroot1974reaching}. This cycling occurs in positive row-stochastic networks without signed edges (antagonistic weights) or adversarial zealots\cite{acemoglu2013opinion,hajnal1958weak}.

The empirical data also reveals a topological immunity condition: kinship clustering\cite{degroot1974reaching}. When a mutually exempt subgraph of agents reaches a critical density threshold ($K \ge N/5$), the network is immunized against ergodicity breakdown\cite{degroot1974reaching,acemoglu2013opinion}. The unpenalized internal connections form an invariant communication core that forces the system into stable convergence or fixed dissensus across all parameter settings\cite{degroot1974reaching,acemoglu2013opinion,hajnal1958weak}.

The remainder of this paper is organized as follows. Section 2 reviews related work across opinion dynamics, bounded confidence, and resilient multi-agent control\cite{degroot1974reaching}. Section 3 formalizes the mathematical preliminaries of discrete consensus and switched affine systems\cite{degroot1974reaching}. Section 4 constructs the status suppression operator, dynamic penalty functions, and kinship clustering metrics\cite{degroot1974reaching}. Section 5 provides theoretical analysis and fixed-point derivations on minimal network realizations\cite{degroot1974reaching}. Section 6 outlines the experimental design, baseline benchmarks, and simulation architecture\cite{degroot1974reaching}. Section 7 presents the empirical results, mapping phase transitions, topological vulnerabilities, and kinship scaling laws\cite{degroot1974reaching}. Section 8 summarizes findings and discusses extensions to continuous-time and multi-issue networks\cite{degroot1974reaching}.

